\section{A secondary continuity clarification}
\label{sec:continuity}
On PDF p.~5, the source attributes its one-fluid density deficit to the
vanishing of $u_y$ at the symmetric mid-plane, describing this as blocking
transport into the sheet~\cite{original}. The value of a flux at one plane
does not determine its divergence:
\begin{equation}
 \partial_t n=-\partial_x(nu_x)-\partial_y(nu_y).
 \label{eq:continuity-law}
\end{equation}
In particular, $\partial_y(nu_y)$ can be nonzero where $u_y=0$.
This is standard continuum kinematics, not a new plasma mechanism. Its role
here is to qualify an explanation based on symmetry alone; the observed
one-fluid density deficit and its dynamical cause remain separate questions.

\begin{proposition}[A periodic density-accumulation example]
\label{prop:stagnation-density}
For $a,n_0>0$, on a periodic coordinate $y\in[0,2\pi)$,
\begin{equation}
 u_y(y)=-a\sin y,\qquad
 n(y,t)=\frac{n_0}{\cosh(at)-\cos y\,\sinh(at)}
 \label{eq:periodic-counterexample}
\end{equation}
is a smooth positive solution of the one-dimensional continuity equation.
Its total mass is $2\pi n_0$, and
$u_y(0)=0$ while $n(0,t)=n_0e^{at}$.
\end{proposition}
\begin{proof}
The denominator is positive for every finite $t$.
Differentiating \eqref{eq:periodic-counterexample} gives
$\partial_tn+\partial_y(nu_y)=0$ directly. The standard integral
$\int_0^{2\pi}(A-B\cos y)^{-1}\dd y=2\pi/(A^2-B^2)^{1/2}$,
with $A=\cosh(at)$ and $B=\sinh(at)$, gives the conserved mass.
Evaluation at zero gives the stated increasing central density.
\end{proof}

This is an exact kinematic transport example, not a self-consistent
plasma--Maxwell solution or a substitute reconnection run. It refutes the
inference from symmetry alone. The two-population model's total density and
density-weighted velocity obey exactly the same continuity equation after
the population equations are summed. A one-fluid model loses the separate
population-crossing information, but it is not geometrically forbidden from
supplying or accumulating total density within a finite-width layer.
Differences in the evolution of aggregate pressure and higher moments can
still explain a particular benchmark failure.
